%------------------------------------------------------------------------------%
\section{Propriedades}

Assim como as funções exponenciais, as funções logarítmicas também possuem uma
lista de propriedades importantes, que decorrem da definição: Sejam
$x, y, z \in \R$ em que $x, y > 0$, então
\begin{gather}
  \log_a (xy)                     = \log_a x + \log_a y \label{eq:log_propriedade_1}\\[3mm]
  \log_a \left(\frac{x}{y}\right) = \log_a x - \log_a y \label{eq:log_propriedade_2}\\[3mm]
  \log_a \left(x^z\right)         = z\log_a x           \label{eq:log_propriedade_3}
\end{gather}
Provaremos a primeira propriedade e as demais serão deixadas como exercício.

%------------------------------------------------------------------------------%
\begin{proof}
  Sejam
  \[
    p = \log_a (xy), \quad \log_a x = m \quad\text{ e } \quad \log_a y = n .
  \]
  Segue da definição de logaritmo que
  \[
    a^p = xy, \quad a^m = x, \quad a^n = y.
  \]
  Portanto substituindo as duas últimas igualdades na primeira obtemos
  \[
    a^p
    = a^m \, a^n
    = a^{m+n}.
  \]
  Logo, $p=m+n$ e verifica-se a propriedade~\eqref{eq:log_propriedade_1}.
\end{proof}
%------------------------------------------------------------------------------%

Mostre as propriedades~\eqref{eq:log_propriedade_2}
e~\eqref{eq:log_propriedade_1} das funções logarítmicas.

%------------------------------------------------------------------------------%
\begin{example}
  Usando o fato de que
  \[
    \log 2 = 0,301, \quad
    \log 3 = 0,477  \quad\text{e}\quad
    \log 5 = 0,699
  \]
  e as propriedades do logaritmo, calcule
  \vspace{-\baselineskip}
  \begin{mathlist}
    \item \log 30
    \item \log 0,333\dots
    \item \log 64
    \item \log \sqrt[7]{2,5}
  \end{mathlist}

  \solution
  \emph{a)}
  \begin{align*}
    \log 30
    & = \log(2 \cdot 3 \cdot 5)  \\
    & = \log 2 + \log 3 + \log 5 & \text{propriedade~\eqref{eq:log_propriedade_1}} \\
    & = 0,301 + 0,477 + 0,699    \\
    & = 1,477
  \end{align*}

  \emph{b)}
  \begin{align*}
    \log 0,333\dots
    & = \log\left(\frac{1}{3}\right) \\
    & = \log 1 - \log 3              & \text{propriedade~\eqref{eq:log_propriedade_2}} \\
    & = 0 - 0,477                    \\
    & =-0,477
  \end{align*}

  \emph{c)}
  \begin{align*}
    \log 64
    & = \log 2^6      \\
    & = 6 \log 2      & \text{propriedade~\eqref{eq:log_propriedade_3}} \\
    & = 6 \cdot 0,301 \\
    & = 1,806
  \end{align*}

  \emph{d)}
  \begin{align*}
    \log \sqrt[7]{2,5}
    & = \log (2,5)^{\sfrac{1}{7}}                  \\[2mm]
    & = \frac{1}{7} \log \left(\frac{5}{2} \right) & \text{propriedade~\eqref{eq:log_propriedade_3}} \\[2mm]
    & = \frac{1}{7} \left(\log 5-\log 2 \right)    & \text{propriedade~\eqref{eq:log_propriedade_2}} \\[2mm]
    & = \frac{0,699 - 0,301}{7}                    \\[2mm]
    & \approx 0,057
  \end{align*}
\end{example}
%------------------------------------------------------------------------------%

%------------------------------------------------------------------------------%
\begin{example}
    Resolva a equação exponencial
    \[
      2^{x-6} = 6,
    \]
    dado que \(\log_3 2 = 1,58\).

    \solution
    Inicialmente, devemos notar que não conseguimos resolver essa equação
    conforme resolvemos algumas na parte de equações exponenciais, uma vez que,
    não conseguimos deixar $6$ na base $2$. Para esses casos, a função
    logarítmica é de extrema importância.

    Sendo assim, vamos aplicar a função $\log_2 x$ em ambos os lados da
    equação. Com isso,
    \begin{align*}
      2^{x-6}                      & = 6                   \\
      \log_2 \left(2^{x-6}\right)  & = \log_2 6            \\
      (x-6) \log_2 2               & = \log_2(2.3)         & \text{propriedade~\eqref{eq:log_propriedade_3}} \\
      x-6                          & = \log_2 2 + \log_2 3 & \text{propriedade~\eqref{eq:log_propriedade_1}} \\
      x                            & = 6 + 1 + 1,58        \\
      x                            & = 8,58
    \end{align*}
\end{example}
%------------------------------------------------------------------------------%

Os dois próximos exemplos mostram que devemos usar a condição de existência do
logaritmo para encontrar o domínio de algumas funções que envolvem logaritmo.

%------------------------------------------------------------------------------%
\begin{example}
  Encontre o domínio da seguinte função
  \[
    f(x) = \log \left( x^2 - 5x + 6 \right)
  \]

  \solution
  A base do logaritmo em questão é 10, com isso, maior que zero e diferente
  de um. Portanto, como não existe logaritmo de números menores ou iguais a
  zero teremos que impor que
  \[
    x^2-5x+6 >0
  \]
  Que é uma inequação do segundo grau cuja solução é
  \[
    S = \left\{ x \in \R \st x<2 \text{ ou } x>3 \right\}.
  \]
  Logo,
  \[
    D(f) = \left\{ x \in \R \st x<2 \text{ ou } x>3 \right\}.
  \]
\end{example}
%------------------------------------------------------------------------------%

%------------------------------------------------------------------------------%
\begin{example}
  Para quais valores de $x$ é possível calcular o logaritmo
  \[
    y = \log_{x+1} \left( -x^2 + 2x + 8 \right)?
  \]

  \solution
  Usando a condição de existência dos logaritmos,
  precisamos impor que
  \begin{gather}
    -x^2 + 2x + 8 > 0 \label{eq:ex_cond_log_1} \\
    x + 1 > 0         \label{eq:ex_cond_log_2} \\
    x + 1 \neq 1      \label{eq:ex_cond_log_3}
  \end{gather}
  Impondo a condição~\eqref{eq:log_propriedade_1}
  \[
    -x^2 + 2x + 8 > 0
  \]
  temos uma inequação do segundo grau, cuja equação associada
  \[
    -x^2 + 2x + 8 = 0
  \]
  possui raízes $x=-2$ e $x=4$.
  O estudo do sinal é apresentado na figura

  \IncludeGraphicCentredEx{parabolalog}[0.5]

  Assim, a fim de que $-x^2+2x+8>0$, devemos ter
  \[
    -2 < x < 4.
  \]
  Além disso, precisamos satisfazer a condição~\eqref{eq:log_propriedade_2}
  \begin{align*}
    x + 2 & > 0 \\
    x     & > -1
  \end{align*}
  e a condição~\eqref{eq:log_propriedade_3}
  \begin{align*}
    x + 1 \neq 1 \\
    x \neq 0
  \end{align*}
  Fazendo a intersecção desses intervalos, temos que
  \[
    y = \log_{x+1}\left(-x^2+2x+8\right)
  \]
  existe no conjunto
  \[
    \left\{ x \in \R \st -1<x<4 \text{ e } x \neq 0 \right\}.
  \]
\end{example}
%------------------------------------------------------------------------------%

%------------------------------------------------------------------------------%
\subsection{Mudança de base}

Sejam $a, b \in (0, 1)\cup(1, \infty)$ e $x\in(0, \infty)$, então
\begin{equation}
  \label{eq:log_mudanca_base}
  \log_{b} x=\frac{\log_a x}{\log_a b}
\end{equation}

%------------------------------------------------------------------------------%
\begin{proof}
  Se
  \[
    z = \log_b x,
  \]
  então
  \[
    b^z=x
  \]
  e
  \[
    \log_a \left(b^z\right) = \log_a x.
  \]
  Mas
  \[
    \log_a \left(b^z\right) = z \log_a b
  \]
  e, portanto
  \[
    z=\frac{\log_a (b^z)}{\log_a b}=\frac{\log_a x}{\log_a b}.
  \]
\end{proof}
%------------------------------------------------------------------------------%

%------------------------------------------------------------------------------%
\begin{example}
    Determine o valor da expressão \(\log_7 25\,\log_5 49\).

    \solution
    Vamos passar $\log_7 25$ para base $5$,
    usando a mudança de base~\eqref{eq:log_mudanca_base}.
    Assim,
    \begin{align*}
      \log_7 25 \,\log_5 49
      & = \frac{\log_5 25}{\log_5 7} \, \log_5 7^2                & \text{usando~\eqref{eq:log_mudanca_base}}  \\[2mm]
      & = \frac{\log_5 5^2}{\cancel{\log_5 7}} 2\cancel{\log_5 7} & \text{usando~\eqref{eq:log_propriedade_3}} \\[2mm]
      & = 4 \log_5 5                                              & \text{usando~\eqref{eq:log_propriedade_3}} \\[2mm]
      & = 4
    \end{align*}
\end{example}
%------------------------------------------------------------------------------%

%------------------------------------------------------------------------------%
\begin{example}
    Escreve a expressão $\log_7 12- \log_7 3$ na base 10.

    \solution
    Vamos primeiro fazer algumas manipulações
    algébricas antes de fazer a mudança de base. Assim,
    \begin{align*}
      \log_7 12 - \log_7 3
      & = \log_7 \frac{12}{3} & \text{usando~\eqref{eq:log_propriedade_3}} \\[2mm]
      & = \log_7 4            & \text{usando~\eqref{eq:log_propriedade_3}} \\[2mm]
      & = \frac{\log 4}{\log 7}
    \end{align*}
\end{example}
%------------------------------------------------------------------------------%

%------------------------------------------------------------------------------%
\subsection{O número de Euler}

Uma base em especial é a mais utilizada em textos científicos, essa base recebe
o nome de \emph{número de Euler}\index{Número de Euler}
\[
  e = 2,718\,281\,82\cdots
\]
As principais funções exponenciais do curso serão $f(x)=e^x$ e $g(x)=e^{-x}$.
O logaritmo com base $e$ recebe o nome de
\emph{logaritmo neperiano}\index{Logaritmo!neperiano} ou
\emph{logaritmo natural}\index{Logaritmo!natural}
e é denotado simplesmente por $\ln$, ou seja,
\[
  \ln x = \log_e x.
\]

%------------------------------------------------------------------------------%